\chapter{Lernziele Kryptologie}

\begin{todolist}
  \item Ich weiss wie die folgenden Kryptosysteme funktionieren (Verschlüsselung und Entschlüsselung bei gegebenem Schlüssel):
  \begin{todolist}
    \item Skytale
    \item Caesar
    \item Vigenère
    \item One-Time-Pad
    \item RSA
  \end{todolist}
  \item Ich kenne die Grundbegriffe: Klartext, Kryptotext, Verschlüsselung (Chiffrierung), Entschlüsselung (Dechiffrierung), Schlüssel, Kryptografie / Kryptologie, Kryptosystem.
  \item Ich kenne die grundlegenden Sicherheitsziele: Vertraulichkeit, Integrität, Authentizität und Verbindlichkeit (Nichtabstreitbarkeit) und kann sie kurz erläutern.
  \item Ich kann das Kerckhoffsche Prinzip erklären.
  \item Ich kann einen Caesar-Kryptotext durch Ausprobieren aller 25 Verschiebungen knacken (Brute-Force) und erkenne den richtigen Klartext.
  \item Ich kann einen Kryptotext, der durch die Vigenère-Verschlüsselung entstanden ist, mit Hilfe des Tools auf folgender Webseite entschlüsseln: \href{https://cryptbreaker.marcwidmer.xyz}{Link zum Cryptbreaker}
  \item Ich kann die Idee der Häufigkeitsanalyse erklären.
  \item Ich kann eine gruppenweise Häufigkeitsanalyse durchführen, um Vigenère bei bekannter Schlüssellänge zu knacken.
  \item Ich kann erläutern, weshalb eine kurze Schlüssellänge Vigenère schwächt (Wiederholungsmuster, Aufteilung in Gruppen) und warum lange bzw. OTP-lange Schlüssel nicht angreifbar durch Häufigkeitsanalyse sind.
  \item Ich kann den Unterschied zwischen monoalphabetischer und polyalphabetischer Substitution erklären.
  \item Ich kann einen monoalphabetisch substituierten Text mittels Häufigkeitsanalyse und schrittweisem Erraten entschlüsseln.
  \item Ich kann die Friedmansche Charakteristik für einen kurzen Text von Hand berechnen.
  \item Ich kann mit der Friedmanschen Charakteristik die Länge eines Vigenère-Schlüssels bestimmen.
  \item Ich kann erklären, was mit der Friedmanschen Charakteristik berechnet wird.
  \item Ich kann beim One-Time-Pad begründen, warum es (bei perfekter Zufälligkeit, einmaliger Verwendung und gleicher Länge wie der Klartext) perfekte Sicherheit bietet (Schlüsselraum = Nachrichtenraum, Gleichverteilung der möglichen Klartexte).
  \item Ich kann die Anzahl möglicher OTP-Schlüssel einer gegebenen Länge berechnen (z.B. $26^n$ für Buchstaben, $2^n$ für Bits).
  \item Ich kann erklären, warum die Wiederverwendung eines OTP-Schlüssels zu Informationsleckage führt und ein Beispiel analysieren.
  \item Ich kann die binäre Version des OTP (Bitweise XOR) anwenden und Kryptotexte berechnen.
  \item Ich kenne die Eigenschaften der XOR-Operation: Kommutativität, Assoziativität, neutrales Element 0, Involution (selbstinverse: $a\oplus a = 0$, $a\oplus 0 = a$) und deren Bedeutung für Ver- und Entschlüsselung.
  \item Ich kann erläutern, warum zweimalige Anwendung desselben Schlüssels (mit XOR) den Klartext zurückgibt.
  \item Ich kann das Three-Pass Protocol (Schlüsseltausch mit drei Durchgängen) Schritt für Schritt durchführen und die beteiligten Nachrichten ($k_A, k_{AB}, k_B$) bestimmen.
  \item Ich kann erklären, warum das Three-Pass Protocol unsicher ist, wenn ein Angreifer alle drei Kryptotexte mitschneidet (Rekonstruktion von $t$ durch XOR aller Nachrichten).
  \item Ich kann das Diffie-Hellman-Merkle-Schlüsseltauschverfahren mit kleinen Zahlen durchführen und den gemeinsamen Schlüssel berechnen.
  \item Ich kann das zugrunde liegende Sicherheitsprinzip von Diffie-Hellman (Schwierigkeit des diskreten Logarithmusproblems) erklären.
  \item Ich kann den Man-in-the-Middle-Angriff auf Diffie-Hellman beschreiben und erläutern, wie digitale Signaturen (z.B. RSA) Authentizität sicherstellen.
  % \item Ich verstehe das Beispiel des \textit{public-key-}Verfahrens mit dem Graphen (Schlüssel: perfect dominating set) und kann erklären, warum hier kein Schlüsselaustausch zwischen Bob und Alice notwendig ist.
  % \item Ich verstehe das Prinzip hinter public-Key-Verfahren.
  \item Ich kann eine Nachricht mit RSA asymmetrisch ver- und entschlüsseln, indem ich ein Beispiel mit kleinen Zahlen durchführe.
  \item Ich kann den Unterschied zwischen symmetrischen und asymmetrischen Verfahren (z.B. Vigenère/OTP vs. RSA/Diffie-Hellman) erklären.
\end{todolist}
